\section{From geometric to energetic epitopes}\label{sec:intro}

\subsection{What the previous method did}
The shipped method (``v2''; it lives at the root of this repository, with \texttt{episcaf\_v3/} as
one subdirectory beside it) scaffolds \emph{B-cell epitopes} --- the regions of an antigen an
antibody recognizes --- so that a high-throughput serology assay (PepSeq) can test
conformational, discontinuous epitopes that linear peptide tiling cannot represent. The epitope is
taken from an antibody:antigen crystal structure as the set of antigen residues the antibody
contacts, defined by a heavy-atom distance cutoff. Those residues are embedded into a \emph{de novo}
protein with RFdiffusion, given a sequence by ProteinMPNN, and folded by AlphaFold3 to check whether
the design actually holds the epitope in its native geometry. Designs are kept or rejected on four
filters: the epitope and overall backbone RMSD to the intended geometry, the predicted alignment
error, and whether scaffold mass clashes into the antibody's binding path. Where an antigen has no
solved antibody, a native-aware cylinder surrogate stands in for that clash filter. A composite
scorer then ranks the survivors and picks the best designs per epitope. This pipeline produced DP4,
a \num{36000}-member PepSeq library now in synthesis.

\subsection{Where the geometric definition runs out}
The method works, but two problems trace to the same root: it treats the epitope as a flat set of
geometric contacts, all equally important.

The first is that binding energy is not uniform across an interface. A handful of ``hot-spot''
residues typically carry most of the binding free energy, surrounded by a ring of residues that
contribute little directly but shape the pocket and exclude solvent. A geometric cutoff cannot see
this: it labels every contacted residue an epitope residue, and the contig we hand RFdiffusion then
fixes the identity and backbone of all of them equally. We are constraining residues we have no
evidence contribute to binding, which costs the generator design freedom it may need to fold a good
scaffold --- and we do not know which contacts are carrying the binding.

The second is visible in the hard cases. A contig enforces epitope residue identity and ordering,
but it does not strongly constrain the distances and orientations \emph{between} spatially separated
epitope islands. So the model can satisfy the contig and still lose the geometry an antibody needs.
The clearest symptom is PDB \texttt{2XQB} patch \texttt{0P}, a 56-residue epitope split across two
islands: across \num{2624} RFdiffusion trials, zero designs passed the filters. A $0\%$ rate over
thousands of attempts is a real failure of what the contig constrains, not undersampling.

Underneath both is a definitional gap. ``Epitope'' is loosely and inconsistently defined in the
literature --- usually by a geometric contact convention chosen for convenience --- and that
ambiguity propagates straight into the design constraints. If we cannot say which residues matter
and why, we cannot say which ones a design is actually required to reproduce.

\subsection{The new avenue: define the epitope by energy}
v3 replaces the geometric definition with an energetic one. Instead of asking which residues
\emph{touch} the antibody, we ask which residues \emph{carry the binding}, and we let the design
constraints follow the energetics. Concretely, we sort the interface into three categories and
constrain each according to what its job is:

\begin{itemize}
  \item \textbf{Category 1 --- energetic (hot-spot) residues.} The residues that dominate the
  binding free energy. These retain both their \emph{exact identity and their shape}: the contig
  fixes their backbone and their sequence is locked to native. This is what we currently do to the
  whole geometric epitope; here we do it only where a residue carries the binding.
  \item \textbf{Category 2 --- structural (occluding) residues.} The rest of the contact shell.
  These matter for how the site fits together --- they hold the shape and exclude solvent around the
  hot spots --- but their identity is not the point. So we fix their \emph{shape} (the contig holds
  the backbone) but let ProteinMPNN choose the sequence. Their correct orientation can be enforced
  by superimposing them on the native antigen and filtering on a small RMSD.
  \item \textbf{Category 3 --- the rest of the scaffold.} These residues exist to enforce the
  geometry of Categories 1 and 2. They can be any identity and adopt any structure, subject to two
  constraints: they must not introduce new steric clashes with the antibody, and they must not
  compete with the binding site. Everything else about them is free --- scaffold length and the
  position of the epitope within the construct are randomized.
\end{itemize}

By fixing identity only where it carries the binding, and shape only where a residue's job is
structural, we give the generator more room to build a foldable scaffold while still holding the
residues that bind. The same energetic reading also sharpens the hard-case problem: it tells us
which inter-residue geometry is worth enforcing, rather than enforcing all of it blindly.

A second change makes the pipeline cheaper and is, in hindsight, overdue: we can filter RFdiffusion
backbones \emph{before} ProteinMPNN and AlphaFold3. A backbone that already places the epitope in
the wrong orientation, or drapes scaffold mass through the antibody's path, was never going to
produce a passing design; superimposing the fixed epitope on the native antigen and checking a small
RMSD plus an antibody-clash test at the backbone stage discards those doomed designs before we spend
the expensive sequence-design and folding steps on them. This is the same accessibility-and-fidelity
check the shipped scorer applies at the end, moved to the front.

\subsection{What this rests on}
The approach depends on getting the energetics right, and the hard part is solvent. A
naive pairwise antigen--antibody interaction energy will mislead: it ignores the desolvation a
residue pays to make a contact, and it misses contributions mediated by ordered water --- a bridging
hydrogen bond can be load-bearing yet invisible to a direct pairwise sum. Deciding how to compute a
defensible per-residue energy, including solvent, is the first real decision of v3, and it is where
the record begins.

We do not take the ranking on trust. The first step is to reproduce known numbers: take a
complex with experimentally measured binding energetics and check that our per-residue energies
recover its known hot spots before we let the definition drive any design. The DP4 binding data,
when it returns, gives a second, independent test --- whether designs that better preserve the
energetic residues bind better.

We are not starting the energy machinery from scratch. The sibling \texttt{bcell\_epitope} project
already computes a per-residue interface interaction-energy decomposition --- including a
water-mediated bridge channel --- on a documented GROMACS stack, and it is the template we build the
energy engine from. We treat it as a reference, not a protocol to inherit: its ``epitopeness'' is an
interaction energy, deliberately not a binding free energy (it omits desolvation and entropy), so
whether v3 adopts that quantity or a fuller one is itself a decision to make and record, not to
assume.

\subsection{Status}
The method above is a plan, not a result: nothing here has been built or tested yet. What follows in this log will be the decisions, one at a time --- the energy method, the
MD protocol, the classification thresholds, the contig strategy, the filters --- each with the
options weighed, the choice made, and the check that backs it, mirrored in \texttt{docs/DECISIONS.md}.
